\documentclass[fontsize=11pt]{article}
\usepackage{amsmath}
\usepackage[utf8]{inputenc}
\usepackage[margin=0.75in]{geometry}
\usepackage{graphicx}
\usepackage{float}
\usepackage{url} 
\usepackage[colorlinks=true, urlcolor=blue, linkcolor=black]{hyperref}

\title{CSC111 Project 2 Report: Modelling and Analyzing the Toronto Urban Public Transit Network using Graph Algorithms}
\author{Shreyas Shastry, Anson Li, Matthew He, Alan Wen Zhuo Xue}
\date{\today}

%Note replace all of the question marks with legit in text MLA citation
\begin{document}
\maketitle
\section*{Introduction}
\textbf{Research question: How does disabling stops in Toronto's public transit network impact overall network efficiency, and which stops are the most essential?}

\vspace{1em}

We've all used the Toronto Transit Commission (TTC), Toronto's public transit network–a massive, interconnected system of subways, streetcars, and buses. Professors, students, and Torontonians all hate TTC disruptions that cause them to miss interviews, midterms, or lectures. To analyze how people travel through the TTC network, we modelled it as a directed, weighted graph using Toronto Open Data datasets \cite{ttc_data}. Vertices represent individual transit stops, and edges represent the directed routes connecting them. We used the normal travel time between stops as the weight for each edge.

\vspace{1em}

Disabling just one major stop can require thousands of commuters to detour, causing a chain reaction that slows down the entire TTC. But from a broader perspective, it is hard to determine which stops are the most critical to keeping the TTC operating efficiently. Our project used two main ideas to explore this issue: computing shortest paths and identifying essential stops. We simulated disruption scenarios computationally by temporarily disabling one stop at a time, by removing the stop and its outgoing edges from the graph and then recomputing travel times. We constrained ourselves to the data given to us by the TTC dataset; when a stop is removed, commuters do not have the option to get off and walk. There must be a valid transit detour when recalculating pathways, or the route is considered "severed" by the disruption.

\vspace{1em}

We reran our shortest-path computations (using the A* algorithm \cite{cormen2022}) for 1,000 sampled origin-destination pairs under each modified network with one stop disabled. To manage the substantial computational runtime of executing this simulation across over 10,000 stops, our program was optimized to utilize memory-efficient subgraph views rather than full graph copies, and we implemented a JSON storing system to save precomputed all-pairs-shortest-paths values on the original graph to speed up subsequent graph traversals.

\vspace{1em}

We defined the impact score of a disabled stop as its network-wide average delay. To calculate this, our algorithm identified the subset of commuters whose original routes relied on the disabled stop and computed their new detour paths. We summed the additional travel time forced upon these affected commuters assigning a strict 120-minute delay penalty if a commuter's route was completely severed and divided this total delay by our entire sample size of 1,000 pairs. The stop that generated the highest network-wide average delay was identified as the most essential stop. Instead of just guessing which stops are the most important, this approach allowed us to objectively identify the network's essential stations. Note that computing the impact scores for 1,000 pairs can be computationally expensive, so our program includes the option to run these calculations with a smaller sample size for a faster, one-minute demonstration.

\vspace{1em}

To measure these disruptions, we first created a "control" of network efficiency by calculating the average travel time of our 1,000 sampled commuters under normal, fully operational conditions with no stops disabled. By comparing this "perfect day" control against the delayed travel times caused by a disabled station, we can figure out how disruptions affect commuters in Toronto as a whole by calculating the average increase in travel time. Ultimately, this tool can help city planners identify Toronto's most critical transit bottlenecks and prioritize maintenance to keep the city moving. If desired, it is theoretically possible to scale our algorithm to simulate the exact population of Toronto provided you have a supercomputer or a few weeks to run the program.

\section*{Dataset Description}
Our project utilizes the official TTC Routes and Schedules dataset, which is published on the City of Toronto Open Data website \cite{ttc_data}. This dataset is provided in the General Transit Feed Specification (GTFS) format, an open standard used globally for public transit schedules. To construct our network for the purposes of our project, we wrote a custom parsing algorithm to extract data from two primary files within this dataset: \texttt{stops.txt} and \texttt{stop\_times.txt} specifically, the rest of the files that can be extracted from the dataset were not used for this project. The other 2 datasets we will mention here, precompute\_apsp.json and pre\_computed\_impact\_scores.json, were generated by our program and were derived from manipulating data from the dataset from the the TTC files we mentioned earlier.

\vspace{1em}

The \texttt{stops.txt} file was used to generate the vertices of our directed graph. Each row in this file corresponds to a specific physical transit stop. For each stop, we extracted its unique stop ID, formal name, and exact geographic coordinates (latitude and longitude). The specific title of each column in this file are the following: stop\_id, stop\_code, stop\_name, stop\_desc, stop\_lat, stop\_lon, zone\_id, stop\_url, location \_type, parent\_station, stop\_timezone, wheelchair\_boarding. However, our project only made use of the relevant columns including stop\_id, stop\_name, stop\_lat, and stop\_lon. These coordinates were essential not only for identifying the stops but also for accurately mapping the final network visually in our MapLibre visualization code. This process populated our graph with over 10,000 distinct nodes that would later be used to connect to neighbouring stops to act as edge weights for the graph. 

\vspace{1em}

To generate the edges and their corresponding weights representing the travel time between 2 neighbouring nodes, we processed the much larger \texttt{stop\_times.txt} file, which contains sequential scheduling data for every individual vehicle trip. The column variables included in that file are as follows: trip\_id, arrival\_time, departure\_time, stop\_id, stop\_sequence, stop\_headsign, pickup\_type, drop\_off\_type, shape\_dist\_traveled. We only made use of the variables trip\_id, arrival\_time, stop\_id, and stop\_sequence. We grouped the data by \texttt{trip\_id} and sorted it by \texttt{stop\_sequence} to trace the exact route order of every bus, streetcar, and subway. This allowed us to identify valid, directed connections between consecutive stops and build the edges that connect all the stops together. Using the arrival time data, we are able to calculate the total time it takes to travel between each station by comparing the arrival times, by converting the HH:MM:SS format into an integer of seconds and dividing by 60 to obtain the amount of minutes to travel between station to station

\vspace{1em}

To make things easier for our TA, we will be submitting or making a program generate a few pre-generated json files that contain the data obtained from running the compute\_impact\_scores function. The first file is called precompute\_apsp and this must be generated after running main.py if you choose to test our compute impact scores function by removing the other json file which we will explain later. This file essentially stores a 2D grid using a list of lists of floats. Each list represents 1 stop, and each index within the length of that list represents every stop in the data and the number associated with each index in that list represents the travel time it takes to go from the stop associated with the outer list to the stop inside the inner list index. Essentially, in our program we use this json file (or we generate it first using one of our functions and then use it) acts as a "guide" or what my teammate likes to call it a heuristic. Normally, the Dijkstra algorithm searches blindly, brute forcing through all possible paths to find the shortest one. So if we already have the graph and how the TTC stations look like on the map and we start at point A in the middle, and we wanna go to point B in the top right, the algorithm has no idea that point B is at the top right, so it will search through all the paths blindly until it finds the shortest one which can take a very long time considering how large the TTC system is. With this generated json file, we now have a way to have the program have some sense of direction which we will explain in further detail in the computation section. For this section the purpose of this dataset is simply for the efficiency of calculating the impact scores in a timely manner and has no presence in the visualization.

\vspace{1em}

The other dataset we will be using is the pre\_computed\_impact\_scores dataset which is provided in the dataset folder but can also be generated again if the TA chooses to delete the data and run our program without it. This dataset is used for our visualization and inside it is simply just a mapping of each stations stop id collected from the stops data file to its associated impact score. This specific file is mainly used in our visualization file, where we take all the stop ids in the stops parameter, and extract the associated latitude, longitude, and station name of that stop id and then use that info to help with placing the vertices on the world map of Toronto using its geographic coordinates and adding interactive hover information for the user to view each stations impact score. Note that if the TA wants to test the generation of both of these datasets and ensure they work, we advise them to head into load\_graph.py, scroll down to around line 47 or simply search for the variable num\_pairs and change the number from 1000 to something small like 50 so the program can run faster. They also must delete the json files from the data for the program to generate the new json files from scratch.

% Here is how it works, lets say we are once again in the middle of the TTC system at point A and we wanna go to the top right at point B. In the code, we mapped each station to a corresponding integer that represents its index in the lists inside the json file. 

% This makes it very easy and efficient to get the index of station A and find the list inside the file associated with station A. Then the code will take the index of station B and find the corresponding travel time from station A and station B using both indices. You may ask, "well you just made it compute all shortest paths, so doesn't that mean you are done after this?" No because if even 1 station is disabled, many routes would be "cut off" and we have to find detours to deal with it. Lets assume the shortest path from A to B had a function C


\section*{Computational Overview}

\vspace{1em}

\textbf{Forms of data our project uses and how graphs play a central role}

\vspace{1em}

As described in the Dataset section we use data provided from the txt and json files as mentioned before but how does our code actually use this in practice? First we have the file graph\_builder.py and this file introduces our TransitGraph class with 2 instance attributes, graph which is a directed graph defined in the networkx library and stops which is a dict mapping stop ids to a tuple in the format (latitude, longitude, stop name). In this class, we initialize the attributes and then we have 2 loading methods to load the stops.txt file and the stop\_times.txt file. Because the files are large we opted to make use of the library pandas to read these files more quickly. The load\_stops method simply adds all the stop ids to the stops dictionary along with the tuple data from reading the file as well as adding the same information as a node of the graph which is helpful for the visualization file to easily get the nodes and edges.

\vspace{1em}

The other method load\_stop\_times loads the entire stop\_times files in small chunks and sorts the data so that we can properly go through each node in the graph and grab the arrival time data of each node. Since the data are sorted we can calculate the time it takes to get from the station we are currently at from the previous station by subtracting the arrival time data in seconds after cleaning it from the station we are at, from the arrival time of the previous station to calculate the weight we will assign to each edge in minutes of travel time. The way we can differentiate between different busses is simply looking at the trip id. If the station I am currently at has the same id, the data is still showing the route the bus is going so we should calculate edge values as normal. If the trip id changes, we are now in a  compeletly different bus and we should "reset" and start calculating travel times for the new route. Note that in the file, arrival time and departure time are the same for each row so to save memory we only used one column.

\vspace{1em}

As you can see the foundation of our project is directly associated with the graph concept, we use vertices to represent stations and we use edges to represent routes. We take this even further by using directed graphs and edge weights because in real life, stations typically go in 1 direction not both ways and the edge weights let us store additional data in the edges that is used for the calculation of impact scores for our project. As we will explain in the following subsections this core structure of a graph connecting each node to its neighbours creates a great way for us to visualize the TTC system on an analytical and geographical level as we can easily obtain the necessary info needed, whether that be station paths, routes, or travel times and use that obtained info to do more exciting things in our program.

\vspace{1em}

\textbf{Describe the major the computations your program performs}

Please note that if you delete the json files in the submission, the program will regenerate them, which will take 1-5 minutes depending on your hardware. After this is complete, you can access the main menu and choose options 1 or 2 for visualizations. The visualization features in our program should appear after a few seconds automatically in your browser. Refresh the page if it's stuck loading.

\vspace{1em}

The important methods for building the graph have already been described. We built the graph using the stops.txt file to represent the nodes an the stop\_times.txt file to represent edges, we also performed data filtering for the edges as a defensive choice in case the times file is corrupted somehow. We also have the \_time\_to\_seconds function to transform the HH:MM:SS strings into integer seconds so we can perform the travel time calculations.

% Matthew here is where you can start going into the Djikstra algorithm stuff I will start it off edit my response as you wish

We designed custom A* and Dijkstra's algorithm implementations in shortest\_path.py to suit the needs of our dataset. Here's a sample run of the computations, assuming that there are no precomputed results stored in .json files. \newline

We start in the load\_graph.py file. First, sample pairs of origins and destinations, representing randomly-generated trips through the TTC network, are generated. We then call the function compute\_impact\_scores (in shortest\_path.py) with the trips. Then, the essential variables needed to store node and edge data are created. The code doesn't find a json file called precomputed\_apsp.json, so it runs all\_paths\_shortest\_paths. \newline

all\_paths\_shortest\_paths finds the shortest time from every stop to every other stop in the TTC (this is called all-pairs-shortest-path, or APSP). This is done by iterating over all stops and calling dijkstra\_all\_nodes, which is a custom Dijkstra's algorithm implementation that finds the shortest distance from a vertex to every other vertex, for each stop. The results are stored in a 2D list. all\_paths\_shortest\_paths has a time complexity of $O(V) \times O((V + E)\log{V}) = O((V^2 + VE)\log{V}$, where $V$ is the number of vertices and $E$ is the number of edges. In practice, it takes about 30 seconds to run for the TTC, which has roughly $V, E = 10^4$. That's still a significant amount of time, so the result is cached in data/precomputed\_apsp.json.
\newline

After that, we find the time and path from the origin to destination for each sample trip. We have the time from computing the APSP, but not the path. Keeping track of the path for each pair of vertices would add more overhead to dijkstra\_all\_nodes than recomputing for each trip, so we recompute. \newline

We use the function a\_star, a custom implementation of the A* heuristic shortest-path algorithm that we created, which returns the shortest distance and path between a source and destination. a\_star uses the APSP values as a heuristic: instead of processing the vertex with shortest distance from the source (as in Dijkstra), it looks at the stop with shortest distance \textit{plus APSP distance from the node to destination}. The A* heuristic for a vertex is an estimation of the cost to get to the destination from that node. For the heuristic to be valid, it must not be greater than the true cost. For the unmodified graph, this means that a\_star will always process the next vertex on the shortest path from the source to destination, with an average running time much faster than regular Dijkstra. This step takes a fraction of a second for $10^3$ trip pairs.

Then, we iterate through each stop, remove it from the network, and recalculate the travel time of each trip to measure the disruption caused. This is the most computationally intensive part of the project. In the previous step, we use the path for each trip to keep track of which trips use each stop (a trip that doesn't use a vertex is not disrupted). For each trip disrupted by a stop removal, we use the a\_star function again to recompute the travel time on the TTC graph without that stop. To avoid $O(n)$ removal of the stop from the graph, the vertex is passed to a\_star's "removed" argument, which is used to avoid relaxing edges to a removed stop. a\_star still works because the heuristic remains valid–removing a node can only increase travel time, so the heuristic for each vertex-destination is never greater than the actual travel time. We calculate the average travel time increases for all trips and store the result in the data folder. \newline 

Using the A* algorithm with APSP on original graph heuristic is about 20x faster than Dijkstra's algorithm, which we used originally. It reduces the processing time for 1000 trips from about ~60 minutes to less than 5 minutes.

\vspace{1em}

\textbf{Results of your Computation in a Visual and/or Interactive way}

Now that we have the graph set up and the path algorithms, how did we display them nicely in a fun way. In main.py after the data is set up, it brings you to a menu allowing you to do 4 things but only the first 2 provide maps, and both of those maps are generated with features from Plotly. Imagine a typical graph, you would think of it as a bunch of circles on a piece of paper with some data inside, and lines connecting circles together. We took some inspiration from this but instead of on paper, we used plotly to display the graph on the map of Toronto. Essentially we gave the necessary plotly functions the geographic data latitude and longitude so it knows where to place the visualizations and where to place each station, and we also used the edge data to tell Plotly which stations need to be connected by an edge to represent a valid route from one station to another. 

\vspace{1em}

We will first describe the first option for our visualization the option to view the entire TTC network map and their impact scores. Using the algorithms and compute impact score functions, every station you see in this feature comes from the impact scores json file that was generated by our programs functions. First we used a feature in Plotly that displays a legend on the right that has a colour map that associates colour with impact score, so if station A has a certain impact score it should have the same colour associated with that score on the legend. Then we took those impact scores that we computed for each station and multiplied that number by 25 and added 5, this was just an arbitrary choice we made to make stations bigger depending on their impact score without making the circles too big. We than used cyan lines to connect the stations represented by dots together to represent the route network of the TTC in a clear visual way that can easily be seen. You can zoom in and out in Toronto as you please and by extracting data such as station name, stop ID and impact score, we made it so the user can view any station they wish by hovering over a station of interest and it will show you the name of the station, the ID of the station, and its given impact score. We did the same for the edges, by extracting the edge weights and names of the connected stations to show the user how long it takes to go from one station to the other based on the data. Note that we chose to assign an invisible vertex in the middle of each edge for edge data to display to avoid having issues with precisely hovering over a certain edge is a lot of edges are around an edge of interest.

\vspace{1em}

The next feature we have is simulating a detour. To use this as the TA you may either find a legitimate route yourself which I personally do not recommend as this is why we chose to implement a feature to do this for you which is option 3, or you can use the 3rd option in our terminal menu and have our program find an example for you to test. When you run this correctly and what you gave as inputs are valid, you will be redirected to the same map and the terminal will output the comparison of travel times between the 2 routes. In the browser you will see 2 routes, a blue route and a red route. The blue route represents the route that a commuter would take in ideal conditions where its the shortest route and there is no disabled stop. The red route is the shortest detour the commuter can take in the event the disabled stop the user inputs is actually disabled which can easily be calculated using the a\_star algorithm that was implemented to find the shortest route quickly even if one stop is missing. As usual we display the node and edge data when you hover over each station or route and the user can see the travel time, station names, etc. We also explicitly labelled the start, destination, and disabled stop in green, yellow, and black explicitly so the user can easily see the stops that they inputted into the terminal.

\vspace{1em}

\textbf{How we use New Libraries to Accomplish Certain Tasks}

\vspace{1em}

We will first discuss pandas.\cite{pandas_readcsv} We used this library in the graph\_builder.py file. This library was not strictly needed but this is a known library to be highly optimized for fast processing of data files and our data was extremely large. Pandas is optimized in C and is generally much faster than python and less straining on your system and it also has some data cleaning features. We used the read\_csv method by pandas to only extract certain columns as shown in the graph builder files since we only need a couple and we used the dtype parameter to ensure that the data is being extracted as the correct data type. We used the method of chunking using chunksize as a paremeter to load the data in batches to prevent the code from trying to do everything at once possibly destroying someones RAM if we process data as large as the stop\_times.txt file. We also used the itertuples method so each row can be converted to a tuple where you can access each item by the name of item you want. Essentially pandas helped us process the data efficiently and saves the user time.

\vspace{1em}

We will also describe networkx. We used it across basically all of our program files as the basis for the directed graph. As said before we mainly use networkx's DiGraph data type\cite{networkx_digraph} which represents a directed graph meaning edges have a single direction they can travel from representing how TTC vehicles travel on a route in 1 direction, and the edges also have weight values that we took advantage of to display travel times between 2 stations. Because this model represents exactly how the TTC system operates on a top down view, we used this datatype to represent the core structure of our graph where nodes represent stations and edges represent route connections and travel times. By the way DiGraph works, when we use the add edge method it takes into the consideration of the order of the parameters to assign the one way direction, in which the first parameter is the origin and the second one is the destination and by this order, the graph will make the original stop have the destination as a neighbour indicating that the station can go there while the destination does not have the origin as a neighbour because it cannot go backwards. We also used the methods has\_path, and subgraph to determine if 2 stations are connected and to assign a subgraph of the original graph as well as shortest\_path and shortest\_path\_length to use the built in Dijkstra algorithms the TA can use to compare. Lastly we used the typical graph features of adding and receiving data such as add\_node, add\_edge, get\_edge\_data, etc. as essential ways for us to build the graph and extract data to be used for stuff like visualization and computing shortest paths.

\vspace{1em}

Lastly we used plotly for the main part of our program which was taking the data we obtained from our other functions and algorithms, and displaying the graph nicely on a world view map of Toronto. We used it in visualization.py in both the create\_transit\_map and show\_shortest\_path functions. The Figure method is like a container or a canvas to draw your map in, it holds all of the important data inside. add\_trace is similar but it adds a sub layer above an existing layer which in this case is the fig variable in the code. On this sub layer we run the go.Scattermap method\cite{plotly_scattermap_api} and this method can take latitude and longitude coordinates and accurately place a dot at that specific geographical location by going through the latitude and longitude coordinate lists at the same time. It can also draw edges, by going through the same process at the vertices but with a None element every 3 elements in the list so the library knows to reset and connect the next 2 stations. Inside this method there are many parameters you can change, such as line or vertex colour, width or size, whether to include hovertext etc.\cite{plotly_scattermap_reference} theres too much to go over in this report but we took advantage of all those parameters to make our visualization colourful and convey information in an appealing way. Lastly we made use of update\_layout to provide the open street map background for our visualization which gives a Google maps kind of map view as the background for the visualization. We use the method fig.show() to display the graph once the user requests to review it in the main terminal in main.py.

\section*{Instructions for Obtaining Datasets and Running the Program}

For our final submission, we have provided a zip file containing the full project directory which. This zip file includes all Python source files, the \texttt{requirements.txt} file, and the \texttt{data/} folder used by the submitted program which will contain all of the files that is extracted from the dataset we mentioned in the dataset section. As a result, the TA does not need to download any datasets from an external URL in order to run the project though if they wish they can refer to the source link in the references section. After extracting the zip folder, all required files should already be present in the same project folder. We will also submit the program files on markus but we cannot upload our data because it is too large. We will provide 2 ways for our TA to download the data, we will email the course email csc111-2026-01@cs.toronto.edu our zip file containing our dataset AND we will also share it through the onedrive option and provide a link as follows when you click it press the CSCProject2-FinalSubmission folder, then download the data folder to get data needed for our project. If the TA downloads the zip folder entirely, the only folder you need is the one called data, please copy paste the entire folder (data folder MUST be included not just individual files in data) and put that folder within the folder you will be using to run our program. The folder structure should essentially be at the outer level, you have your folder storing our entire project, and inside that folder you should have all of our markus submitted files as well as the data folder inside the outer folder, and when opening the data folder you should see the dataset txt files along with a json file.

\url{https://utoronto-my.sharepoint.com/:u:/g/personal/matt_he_mail_utoronto_ca/IQA7-vWOdwgxQIOYF4j3JntfAbFnJlshT0b-BnMB6YLmfFg?e=fsTT6y} \newline

We also include precomputed delay data for 1000 random trips, but not the all-pairs-shortest-path (APSP) data because it's too large. To test the APSP feature and the delay calculation feature, delete the precomputed json files in the data/ directory. It runs in less than 5 minutes. If it takes too long, change the "num\_pairs" variable in load\_graph.py to 100 instead of 1000.

\vspace{1em}

In particular, the extracted project folder should contain:
\begin{itemize}
    \item \texttt{main.py}
    \item the supporting Python modules such as \texttt{graph\_builder.py}, \texttt{load\_graph.py}, \texttt{shortest\_path.py}, \texttt{visualization.py}, and any other submitted project files
    \item \texttt{requirements.txt}
    \item The dataset zip file (separate) contains a \texttt{data/} directory containing the files needed by the main program, including \texttt{stops.txt}, \texttt{stop\_times.txt}, and \texttt{precomputed\_impact\_scores.json}
\end{itemize}

\vspace{1em}

To run the project, the TA should:
\begin{enumerate}
    \item Install Python 3.13.5.
    \item Extract the submitted project zip into a folder.
    \item Open that folder in PyCharm or open a terminal in that folder.
    \item Download and unzip the dataset zip.
    \item Move the data/ directory in the dataset folder into the main project directory.
    \item Install all required external libraries by running \texttt{pip install -r requirements.txt}.
    \item Make sure the \texttt{data/} directory remains in the same folder as \texttt{main.py}.
    \item Run \texttt{main.py}.
\end{enumerate}

\vspace{1em}

When main.py is ran by the TA, the TA should see the code output loading transit data and then processing chunk lines that print out every couple of seconds and it should look something like this below after a few seconds assuming that the TA is NOT trying to generate the json files we gave themselves and already have the ones we provided in the zip folder. All of the output at this point is strictly inside the python terminal.

\begin{figure}[H]
    \centering
    \includegraphics[width=1.0\linewidth]{Main menu terminal screen.png}
    \caption{Main Menu view after a few seconds}
    \label{fig:placeholder}
\end{figure}

\vspace{1em}

\newpage
\textbf{Option 1: View The Network Impact Score Map.}  
This opens an interactive map in the browser showing TTC stops and their impact scores along with the edges connecting each station to its neighbours. This should only take a few seconds to automatically pop up in your browser, sometimes you may need to refresh the browser tab that it brings you to if it starts infinitely reloading. Stops with higher impact scores appear larger and scale with the colour legend on the right. Hovering over stops reveals the stop name, stop ID, and impact value. You may also click trace on the top right to remove the stations from your view. Hovering over route segments can also show travel-time information for edges. The TA should be led to a page that looks like this: 

\begin{figure}[H]
    \centering
    \includegraphics[width=1\linewidth]{Option 1.png}
    \caption{Impact Score Transit Map}
    \label{fig:placeholder}
\end{figure}

\vspace{1em}

\newpage
\textbf{Option 2: Simulate a Custom Route Detour.}  
The user enters an origin stop ID, a destination stop ID, and a stop ID to disable. The program first checks whether a shortest path exists in the intact network. If the disabled stop lies on the original shortest path, the program computes a detour route in a temporary subgraph with that stop removed if using the built-in djikstra algorithm. If the TA wishes to use our custom algorithm, they may call themselves and pass in the correct parameters. It then prints the original and new travel times and opens a second map showing the original route, the detour route, the start and end points, and the disabled stop. You may need to zoom out to see the routes depending on where the stops are. In the terminal it will also display the travel time using the normal route  and the detour route as a comparison. It should look something like the following image for the visualization:

\begin{figure}[H]
    \centering
    \includegraphics[width=1\linewidth]{Option 2.png}
    \caption{Detour Map Visualization}
    \label{fig:placeholder}
\end{figure}

\vspace{1em}

\textbf{Option 3: Auto-Generate Valid Detour Scenario.}  
This option searches for a pair of stops and a removable stop that produce a noticeable delay. It prints example stop IDs that can then be copied into Option 2. This is useful for quickly demonstrating the detour functionality without needing to manually search the full dataset yourself which we tried and was not fun. The terminal will give an output like this:

\begin{figure}[H]
    \centering
    \includegraphics[width=1\linewidth]{Option 3.png}
    \caption{Terminal Displaying Valid Example}
    \label{fig:placeholder}
\end{figure}

\vspace{1em}

\textbf{Option 4: Quit.}  
This exits the program. An image should not be required here it simply just stops main.py.

\vspace{1em}

\section*{Description of Changes}

Our final submission remains very close to the plan proposed in Phase 1, but a couple important implementation details changed as we moved from a proposal to a working program.

\vspace{1em}

First, we mentioned in the proposal that we would use directed graphs but we did not specify that were going to use the one provided by networkx. Given that we wanted to implement our own path finding algorithm, also creating a directed graph class from scratch would make it quite time consuming to do both complex implementations from scratch. Given that the pathfinding algorithm was something we were interested in, we opted to focus on that and make use of the DiGraph class from the library. Nonetheless, we were still able to fulfill the report expectations for this feature as we still used the directed graphs for our visualizations, our algorithms etc and we demonstrated our proficiency with the new library as well.

\vspace{1em}

Another thing was we considered our TA's feedback about managing runtime. Although their feedback was mainly focused on the proposal, they mentioned that we should optimize on running time for the path finding algorithm. Initially we planned on generating maybe a 50 sample pair json file for the TA to run because it was very slow at the start, but after we continued optimizing the algorithm it came to the point where even the 1000 pair runtime was becoming quite efficient. This was not explicitly planned in the proposal but it came to fruition as we started focusing on it significantly improving the efficiency of the algorithm.

\vspace{1em}

Finally, the final program became more interactive than our original proposal described. In addition to producing an overall network impact score map, the submitted code includes a custom detour simulator inside main.py. These features help visualize the Dijkstra algorithm in practice and how it effectively deals with station disruptions in the shortest path it can find. 

\section*{Discussion}

\textbf{Do the results of your computational exploration help answer this question/achieve your goal?}

\vspace{1em}

Our projects computational algorithms and visualizations significantly helped us calculate and answer our research question. As we disable stops in the TTC network, it is evident by our visualization that many stations in the TTC system affects the overall network in at least a moderate way as evident by the amount of stations you can see on the map that display as a colour that has a medium to high impact score. The use of Dijkstra's algorithm to compute the data needed and using Plotly's libraries to display it nicely has helped us conclude which stops in the TTC are most essential (which in the provided file case, would be the Exhibition Loop, Queensway and the stations around Bathurst Street as shown in the image below) and it also verifies that disabling stops in the TTC has varying impact on the general network but in general, the impact is marginally low.

\begin{figure}[H]
    \centering
    \includegraphics[width=0.8\linewidth]{Example.png}
    \caption{The bright orange/yellow vertices have highest impact scores}
    \label{fig:placeholder}
\end{figure}

One important insight from the results of this project is that a stop's practical importance depends on network structure, not just how recognizable the stop is to riders. A good example we thought of on this map is Union Station, from our experience in Toronto Union station is by far what we would think as a very busy station and what we would expect to have a significant impact on the network, but as you will see below, if Union is somehow disabled, the impact is still moderate in comparison to other stations on the map. On the other hand, a stop that is less famous may become extremely important if it acts as a narrow bridge between parts of the network. This was the case when we analyzed the yellow vertex on the very bottom of the map in figure 5, at least we have not heard much about this stop and yet it has one if not the highest impact on the network.

\begin{figure}[H]
    \centering
    \includegraphics[width=1\linewidth]{Example2.png}
    \caption{Union Stations score is only 0.84 which is less than half the score of the Exhibition Loop}
    \label{fig:placeholder}
\end{figure}

\vspace{1em}

Another strength of the final program is that it works at multiple levels of analysis without being too complex for the user. They could run the first option to learn about the TTC and what the network looks like, and which stations need extra attention from the city to reduce network disruptions and they could theoretically also use our program to find an alternate detour for them to take if they want to go from one place to another but one of the stations is closed for whatever reason. Not only does our program produce results that answer our research question, it can also be easily used and can provide actual benefits for users if they find that their normal route is blocked off by a disabled stop (and they refuse to use existing apps like Google Maps).

\newpage
\textbf{What limitations did you encounter, with the datasets you found, the algorithms/libraries you used, or other obstacles?}

\vspace{1em}

One limitation that our proposal TA mentioned is that the data we are using is based on schedule data and not live data for real time operation. We are simply using pre-collected data provided by the TTC and generating a static view of a snapshot in time of the data we had for this project. 

This leads into another limitation with the dataset in which this dataset essentially only considers routes in the TTC, it may be the case that it would be faster for a commuter to get off the bus at some point, walk to the next line, and get on that bus to get to their destination faster, although realistically people would do that, it would be very complex to implement because we would have to consider the average walking speed, and would it actually be that much faster, and we would need to connect additional edges based on distance and not off a dataset.

A particular limitation we ran into for our libraries was the speed of the Dijkstra algorithm built by networkx. We used this method mainly when we did not have the custom one built yet, and although it works great with the custom route detour feature, it is incredibly slow at finding the shortest path when we tried to use it with compute\_impact\_scores it would take almost 2 hours for it to run and give us the json file. In fact, our custom algorithm ended being much faster at doing the exact same task so we ended up replacing it entirely with our custom made algorithm to do the computing.

\vspace{1em}

\textbf{What are some next steps for further exploration?}

\vspace{1em}

There are a few next steps for future work on this project. One improvement would be to add walking-transfer edges between nearby stops so that the model better reflects real commuter behaviour and although this is unlikely to be entirely perfect and accurate, it would likely be a more realistic approach when considering the options commuters have to get to their destination. We could also distinguish between buses, streetcars, and subways more explicitly and display them as such in our visualization or incorporate time-of-day information, since a stop's importance may vary between peak and off-peak periods as transit becomes important when people leave for and from work.

\vspace{1em}

Overall, the project succeeds in demonstrating that graph algorithms provide a useful structure for studying transit disruptions and can easily be visualized with the right tools and libraries. Even with the limitations of schedule data and sampled shortest paths, the project provides a concrete answer with substantial evidence backing it up to the question of which TTC stops are most essential and the significance of the impact on the TTC when stations are disabled. It also produces results in a form that is easy to understand through both an overall map and specific route-level detours, making the analysis more understandable to a user than simply a data table that is boring to read.



\newpage

\begin{thebibliography}{9}

    \bibitem{cormen2022}
    Cormen, Thomas H., Charles E. Leiserson, Ronald L. Rivest, and Clifford Stein.
    \textit{Introduction to Algorithms}. 4th ed., MIT Press, 2022.

    \bibitem{networkx_digraph}
    NetworkX Developers.
    ``DiGraph Directed graphs with self loops.''
    \textit{NetworkX Documentation},
    
    \url{https://networkx.org/documentation/stable/reference/classes/digraph.html}.
    Accessed 31 Mar. 2026.

    \bibitem{pandas_readcsv}
    Pandas Development Team.
    ``pandas.read\_csv.''
    \textit{pandas Documentation},
    
    \url{https://pandas.pydata.org/docs/reference/api/pandas.read_csv.html}.
    Accessed 31 Mar. 2026.

    \bibitem{plotly_scattermap_api}
    Plotly.
    ``plotly.graph\_objects.Scattermap.''
    \textit{Plotly Python API Reference},
    
    \url{https://plotly.com/python-api-reference/generated/plotly.graph_objects.Scattermap.html}.
    
    Accessed 31 Mar. 2026.

    \bibitem{plotly_scattermap_reference}
    Plotly.
    ``Python Figure Reference: scattermap Traces.''
    \textit{Plotly Graphing Libraries},
    
    \url{https://plotly.com/python/reference/scattermap/}.
    Accessed 31 Mar. 2026.

    \bibitem{ttc_data}
    Toronto Transit Commission.
    ``TTC Routes and Schedules.''
    \textit{City of Toronto Open Data Portal}, 3 Feb. 2026,
    \url{https://open.toronto.ca/dataset/ttc-routes-and-schedules/}.
    Accessed 31 Mar. 2026.

\end{thebibliography}

\end{document}




